\section{Locate rows in memory}\label{blit-x-layout}
\begin{table}[htbp]
  \centering
  \caption{Default row geometry. All byte offsets are relative to the selected allocation base.}
  \tagpdfsetup{table/header-rows={1}}
  \begin{tabular}{lrrr}
    \toprule
    Object      & Width & Word stride & Row-one byte offset \\
    \midrule
    Framebuffer & 8     & 16          & 64                  \\
    Sprite      & 5     & 7           & 28                  \\
    \bottomrule
  \end{tabular}
\end{table}
A copy width of three does not make either row stride twelve bytes. After the
inner loop consumes its three source cells, the next source row still begins
28 bytes after the previous row start. Likewise, the next destination row
begins 64 bytes after its previous row start. The guided implementation retains
these starts instead of reconstructing them from the advanced cursors.

For a base subview beginning at word offset three, add twelve bytes to the
allocation base before applying row and column offsets. The descriptor already
contains this adjusted address. Adding that displacement again in the function
would select the wrong cells. The public checks choose such subviews independently
for source and destination.
